\documentclass[11pt,a4paper]{article}

\usepackage[utf8]{inputenc}
\usepackage[T1]{fontenc}
\usepackage{lmodern}
\usepackage[french]{babel}
\usepackage[margin=2.3cm]{geometry}
\usepackage{microtype}
\usepackage{booktabs}
\usepackage{array}
\usepackage[table]{xcolor}
\usepackage{amsmath}
\usepackage{amssymb}
\usepackage{mathtools}
\usepackage{enumitem}
\usepackage[most]{tcolorbox}
\usepackage{tabularx}
\usepackage{caption}
\usepackage{titlesec}
\usepackage{hyperref}

\definecolor{navy}{HTML}{1F3864}
\definecolor{bluemid}{HTML}{2E5496}
\definecolor{lightblue}{HTML}{D6E0F0}
\definecolor{accent}{HTML}{C0491F}
\definecolor{green}{HTML}{2E6B4F}
\definecolor{lightgreen}{HTML}{D6EADF}

\hypersetup{colorlinks=true, linkcolor=bluemid, urlcolor=bluemid,
  pdftitle={Ossature SCR - construction chiffree DORA},
  pdfauthor={Kelian Kaddouri}}

\newtcolorbox{cle}[1][]{colback=lightblue!40,colframe=navy,boxrule=0.6pt,
  arc=2pt,left=8pt,right=8pt,top=6pt,bottom=6pt,before skip=9pt,after skip=9pt,#1}
\newtcolorbox{attention}[1][]{colback=accent!8,colframe=accent,boxrule=0.6pt,
  arc=2pt,left=8pt,right=8pt,top=6pt,bottom=6pt,before skip=9pt,after skip=9pt,#1}
\newtcolorbox{apport}[1][]{colback=lightgreen!60,colframe=green,boxrule=0.6pt,
  arc=2pt,left=8pt,right=8pt,top=6pt,bottom=6pt,before skip=9pt,after skip=9pt,#1}

\titleformat{\section}{\normalfont\large\bfseries\color{navy}}{\thesection}{0.6em}{}
\titlespacing*{\section}{0pt}{2.0ex plus .3ex minus .2ex}{1.0ex plus .2ex}
\titleformat{\subsection}{\normalfont\normalsize\bfseries\color{bluemid}}{\thesubsection}{0.6em}{}
\captionsetup{font=small,labelfont={bf,color=navy},labelsep=period,skip=6pt}

\newcommand{\Ind}{\mathbf{1}}
\newcommand{\E}{\mathbb{E}}
\newcommand{\Var}{\mathrm{Var}}
\newcommand{\VaR}{\mathrm{VaR}}
\newcommand{\TVaR}{\mathrm{TVaR}}

\title{\vspace{-1.2cm}\color{navy}\bfseries Ossature du SCR DORA\\[2pt]
  \large De la cascade qualitative a une charge de capital : conventions,
  severite, frequence, agregation, mesure de risque}
\author{Kelian Kaddouri}
\date{\today}

\begin{document}
\maketitle
\vspace{-0.5cm}

\begin{cle}
\textbf{Ce que fait ce document.} Il pose le squelette mathematique de la partie
chiffree. La brique \emph{severite} (section~3) est entierement specifiee, car
c'est celle qu'implemente le script \texttt{07\_severite.py}. Les briques
\emph{frequence}, \emph{agregation} et \emph{mesure de risque} (sections~4 a~6)
sont posees comme charpente, a remplir par les scripts \texttt{08} a \texttt{10}.
\end{cle}

\begin{attention}
\textbf{Cadre honnete, a garder present partout.} DORA (UE 2022/2554) impose de la
resilience operationnelle, \emph{pas} une exigence de capital : il n'existe pas de
SCR DORA reglementaire. On construit ici cet objet par analogie avec Solvabilite~II.
Sans donnees de perte calibrees sur l'entite, le \emph{niveau absolu} du SCR n'est
pas identifiable : seules sa \emph{structure} et sa \emph{sensibilite} le sont. On
travaille donc en unites normalisees et le livrable est la surface $\mathrm{SCR}(g,\xi)$,
pas un montant en euros.
\end{attention}

% ============================================================
\section{Conventions de modelisation}

\begin{apport}
\textbf{Deux echelles de dependance, disjointes.} C'est le choix structurant qui
evite le double comptage.
\begin{itemize}[leftmargin=1.4em,itemsep=1pt,topsep=2pt]
  \item \textbf{Inter-incidents} : un facteur systemique commun $Y\sim\mathcal N(0,1)$
        gouverne la \emph{frequence} des incidents et la co-occurrence des cellules
        (une mauvaise annee en declenche beaucoup).
  \item \textbf{Intra-incident} : la cascade dirigee (\texttt{ROOT}/\texttt{TRANS})
        gouverne l'\emph{etendue} d'un incident donne, l'ensemble $S$ des piliers
        touches. Le gain de propagation $g$ est deterministe en cas de base ($g$
        pilote par $Y$ est mentionne comme extension, non calibre).
  \item \textbf{Severites} tirees \emph{independamment sachant $S$} : la dependance
        est deja portee par la formation de $S$ (intra) et par $Y$ (inter). Aucun
        troisieme canal.
\end{itemize}
\end{apport}

% ============================================================
\section{Notations}

$P=\{1,\dots,5\}$ les cinq piliers DORA ; $\mathrm{GBASE}_j$ la gravite de base
ordinale du pilier $j$ ; $S\subseteq P$ l'ensemble des piliers touches par un
incident (issu du moteur de cascade auto-evitant) ; $L_j$ la perte generee par le
pilier $j$ ; $X(S)=\sum_{j\in S} L_j$ la severite d'un incident ; $N_{ij}$ le nombre
d'incidents de la cellule $(i,j)$ sur un an ; $\Phi$ la fonction de repartition
normale standard, $\phi$ sa densite.

% ============================================================
\section{Severite : de l'ordinal aux unites monetaires}

\subsection{Loi de perte d'un pilier}

Le score $\mathrm{GBASE}_j$ ne fixe pas un montant, il fixe l'\emph{echelle} de la
loi du pilier $j$. La forme est un corps lognormal prolonge par une queue de Pareto
generalisee (GPD) au-dela d'un seuil (methode des exces, POT). L'indice de queue
$\xi$ est \textbf{commun} a tous les piliers : c'est un regime de queue du risque
cyber, pas une propriete d'un pilier.

\begin{equation}
\ln L_j \sim \mathcal N(\mu_j,\sigma^2)\ \text{dans le corps},\qquad
\mu_j=\mu_0+c\,(\mathrm{GBASE}_j-\min_k \mathrm{GBASE}_k).
\end{equation}

Le seuil de raccord $u_j$ est le quantile de niveau $p_u$ du corps, et l'echelle de
queue $\beta_j$ est \textbf{derivee} par continuite de la densite au raccord (elle
n'est pas un parametre libre) :
\begin{equation}
u_j=\exp\!\big(\mu_j+\sigma z_{p_u}\big),\quad z_{p_u}=\Phi^{-1}(p_u),\qquad
\beta_j=\frac{1-p_u}{f_{\mathrm{LN}}(u_j)}=\frac{(1-p_u)\,u_j\,\sigma}{\phi(z_{p_u})}.
\end{equation}

Au-dessus du seuil, la survie suit la GPD :
\begin{equation}
\Pr(L_j>x)=(1-p_u)\Big[1+\xi\,\tfrac{x-u_j}{\beta_j}\Big]^{-1/\xi},\qquad x>u_j .
\end{equation}

\begin{attention}
\textbf{Ce qui est pose, pas estime.} Le point de raccord $p_u$ et la pente $c$ de la
carte $\mathrm{GBASE}\!\to\!\mu$ sont des hypotheses, non des estimations sur donnees
de l'entite. L'ancrage des unites (le niveau de $\mu_0$) est une convention de
normalisation. On l'assume ; on ne le deguise pas en mesure.
\end{attention}

\subsection{Simulation exacte par inverse de la fonction de repartition}

Pour $U\sim\mathcal U(0,1)$ :
\begin{equation}
L_j=
\begin{cases}
\exp\!\big(\mu_j+\sigma\,\Phi^{-1}(U)\big) & U<p_u \quad(\text{corps}),\\[4pt]
u_j+\dfrac{\beta_j}{\xi}\Big[(1-V)^{-\xi}-1\Big],\ \ V=\dfrac{U-p_u}{1-p_u} & U\ge p_u \quad(\text{queue}).
\end{cases}
\end{equation}

La moyenne de queue est finie car $\xi<1$ : l'exces moyen au-dela du seuil vaut
$\beta_j/(1-\xi)$. Pour $\xi\ge 1$ la moyenne diverge : on impose donc $\xi<1$ et on
signale l'instabilite numerique quand $\xi$ approche $0{,}9$.

\subsection{Severite d'une cascade}

Un incident touche l'ensemble $S$ ; chaque pilier touche tire sa perte, et on somme :
\begin{equation}
X(S)=\sum_{j\in S} L_j .
\end{equation}
Les $L_j$ sont independantes sachant $S$ (conformement aux conventions). C'est le
point de jonction avec la cascade qualitative : le moteur \texttt{ROOT}/\texttt{TRANS}
produit $S$, la loi de severite produit $X(S)$.

% ============================================================
\section{Frequence a facteur systemique commun (charpente, script 08)}

Un unique facteur $Y\sim\mathcal N(0,1)$ pilote la surdispersion et la co-occurrence :
\begin{equation}
N_{ij}\mid Y \sim \mathrm{Poisson}\!\big(\lambda_{ij}\,e^{\,a_j Y-\tfrac12 a_j^2}\big),
\qquad \E[N_{ij}]=\lambda_{ij},\quad \Var(N_{ij})>\E[N_{ij}].
\end{equation}
Le terme $-\tfrac12 a_j^2$ recentre l'esperance ; l'intensite de $Y$ ($a_j$)
remplace une binomiale negative bricolee a cote d'un Vasicek : une seule source de
risque systemique, celle du lab existant.

% ============================================================
\section{Agregation et mesure de risque (charpente, scripts 09 et 10)}

Perte agregee sur un an, par Monte-Carlo (tirer $Y$, puis $N_{ij}$, puis pour chaque
incident son ensemble $S_k$ et sa severite $X_k$) :
\begin{equation}
S_{\text{total}}=\sum_{(i,j)}\ \sum_{k=1}^{N_{ij}} X_k(S_k).
\end{equation}

Le SCR est la Value-at-Risk a $99{,}5\%$, par mandat Solvabilite~II (art.~101) :
\begin{equation}
\mathrm{SCR}=\VaR_{0{,}995}\big(S_{\text{total}}\big).
\end{equation}

\begin{attention}
\textbf{VaR, pas TVaR, et pour la bonne raison.} On retient la VaR parce que la norme
Solvabilite~II l'impose, \emph{pas} parce que la TVaR serait incalculable : la TVaR
reste finie tant que $\xi<1$. On la reporte a titre indicatif
($\TVaR_{0{,}995}=\VaR_{0{,}995}/(1-\xi)+\dots$) sans en faire la mesure retenue.
Le centrage (SCR brut $\VaR$, ou capital economique $\VaR-\E$) est precise au calcul.
\end{attention}

\begin{apport}
\textbf{Le livrable.} Pas un nombre mais la surface $\mathrm{SCR}(g,\xi)$ : le capital
en fonction du gain de propagation $g$ et de l'indice de queue $\xi$, avec la zone
instable ($\xi>0{,}9$) marquee et la frontiere identifiable / non-identifiable tracee
dessus. Le message : le capital est une fonction d'hypotheses expertes explicites, et
voici lesquelles le deplacent le plus.
\end{apport}

\end{document}
